\section{Conventions}
\label{app:conventions}

Everything in this paper depends on a consistent set of conventions, and several
of the results are sensitive to sign choices in a way that is easy to
underestimate. They are collected here in full.

\paragraph{Metric and signature.}
Spacetime is $\mathbb{R}^{1,9}$ with
$\eta_{\mu\nu} = \mathrm{diag}(-1,+1,+1,\dots,+1)$, indices
$\mu = 0,1,\dots,9$. Indices are raised and lowered with $\eta$. The split
signature $(5,5)$ appears only as the natural home of the oscillator frame in
appendix~\ref{app:orientation}; it is a different real form of the same complex
algebra, not a change of frame within the Lorentzian one.

\paragraph{Orientation and epsilon.}
The orientation is fixed by $\epsilon_{01\dots 9} = +1$, hence
$\epsilon^{01\dots 9} = -1$ in Lorentzian signature. The volume form is
$\mathrm{vol} = \tfrac{1}{10!}\,\epsilon_{\mu_1\dots\mu_{10}}\,
 dx^{\mu_1}\wedge\dots\wedge dx^{\mu_{10}}$.

\paragraph{Hodge star and self-duality.}
On five-forms,
\begin{equation}
(\star F)_{\mu_1\dots\mu_5}
 = \frac{1}{5!}\,\epsilon_{\mu_1\dots\mu_5}{}^{\nu_1\dots\nu_5}
   F_{\nu_1\dots\nu_5},
\end{equation}
and with the above conventions $\star^2 = +\starSquared$ on middle-degree forms
in ten Lorentzian dimensions. This sign is what makes the eigenspaces real; it is
verified computationally rather than assumed, and a test asserts it directly.

The sign is signature-dependent: in Euclidean signature $\star^2 = -1$ on
middle-degree forms in ten dimensions, the eigenvalues are $\pm i$, and there is
no real eigenspace decomposition at all. We state this here, where the convention
is fixed, rather than only in passing, because this project once mistook a
Euclidean projector for an anti-self-dual one and read the resulting
non-vanishing as a signal (appendix~\ref{app:gten}). Any routine taking a
signature flag must therefore be called deliberately, and the value of
$\star^2$ is checked rather than presumed wherever the decomposition is used.
A five-form is self-dual when $\star F = +F$ and anti-self-dual when
$\star F = -F$; both eigenspaces have dimension $\dimSelfDualModule$.

\paragraph{Inner product.}
$\langle F, G\rangle = F_{\mu_1\dots\mu_5}G^{\mu_1\dots\mu_5}$, with all five
indices raised by $\eta$. This is the pairing used throughout
appendix~\ref{app:gten}.

\paragraph{Clifford algebra.}
$\{\Gamma_\mu, \Gamma_\nu\} = 2\eta_{\mu\nu}\mathbf{1}$, with
$32\times 32$ gamma matrices. Antisymmetrised products are written
$\Gamma^{\mu_1\dots\mu_k} = \Gamma^{[\mu_1}\cdots\Gamma^{\mu_k]}$ with unit
weight, so that $\Gamma^{\mu_1\dots\mu_k}$ agrees with
$\Gamma^{\mu_1}\cdots\Gamma^{\mu_k}$ when all indices are distinct.

\paragraph{Chirality.}
$\Gamma_{11} = \Gamma^0\Gamma^1\cdots\Gamma^9$ satisfies $\Gamma_{11}^2 = +1$,
and the chiral projectors are $P_\pm = \tfrac12(1\pm\Gamma_{11})$. The chiral
module $S_+$ has dimension $\spinorDim$.

\paragraph{Charge conjugation.}
$C$ satisfies $C\Gamma^\mu C^{-1} = -(\Gamma^\mu)^{T}$ in the convention used
here. The relevant consequence is the symmetry of $C\Gamma^{(5)}$ on the chiral
subspace, which is what allows the bridge \eqref{eq:bridge} to land in
$\mathrm{Sym}^2 S_+$. This symmetry is checked directly in the implementation.

\paragraph{Five-gamma duality.}
On the chiral subspace $\Gamma^{(5)}$ is not independent of its dual: the
self-dual part of $\Gamma^{(5)}$ acts non-trivially and the anti-self-dual part
annihilates $S_+$. This is the representation-theoretic reason the bridge sees
only one eigenspace, and it is the fact that the orientation branch of
appendix~\ref{app:orientation} can silently invert.

\paragraph{Normalisation.}
The $1/5!$ in \eqref{eq:bridge} is a choice. Every dimension, rank and
intersection reported in this paper is invariant under rescaling the bridge, so
nothing depends on it. It is fixed only so that the round-trip identity
$\Phi^-\Phi = \mathrm{id}$ holds on the nose rather than up to a constant.

\paragraph{Field degree.}
``Degree'' always means degree in $F$, never degree as a differential form. An
invariant of field degree $d$ is a polynomial of degree $d$ in the
$\dimSelfDualModule$ components of $F$. Only even field degrees occur.
